\section{From phase pruning to certified SKAT tails}

The previous sections certify the effect of phase pruning on Krylov
propagation.  A SKAT-like test ultimately depends on the spectrum itself, so
one more bridge is needed:
\[
\text{phase pruning}
\longrightarrow
\text{operator error}
\longrightarrow
\text{spectral error}
\longrightarrow
\text{tail-probability error}.
\]

\subsection{Block-norm operator domination}

Let
\[
E=H-H_0
\]
be the self-adjoint error introduced by phase pruning.  Using the same
orthogonal phase projectors \(P_h\), define
\[
\mathcal C_{hg}
=
\|P_hEP_g\|_2.
\]

\begin{theorem}[Block-norm operator certificate]
\[
\boxed{
\|E\|_2
\le
\|\mathcal C\|_2.
}
\]
\end{theorem}

\begin{proof}
For any vector \(x\), let
\[
r_g=\|P_gx\|_2.
\]
Then
\[
\|P_hEx\|_2
\le
\sum_g
\|P_hEP_g\|_2\,\|P_gx\|_2
=
(\mathcal C r)_h.
\]
Orthogonality of the phase decomposition gives
\[
\|Ex\|_2^2
=
\sum_h\|P_hEx\|_2^2
\le
\|\mathcal C r\|_2^2
\le
\|\mathcal C\|_2^2\|r\|_2^2
=
\|\mathcal C\|_2^2\|x\|_2^2.
\]
Taking the supremum over unit vectors proves the claim.
\end{proof}

The important computational point is that \(\mathcal C\) is only a
phase-by-phase matrix.  The same block norms already used to construct the
polyphase scheduler therefore provide a global operator certificate.

\subsection{Weyl enclosure of the quadratic-form coefficients}

Assume
\[
H\succeq0
\]
and let the eigenvalues of \(H\) and \(H_0\), in decreasing order, be
\[
\lambda_i
\qquad\text{and}\qquad
\mu_i.
\]
If
\[
\|H-H_0\|_2\le\varepsilon,
\]
Weyl's inequality gives
\[
|\lambda_i-\mu_i|\le\varepsilon.
\]

\begin{corollary}[Positive spectral interval]
For every \(i\),
\[
\boxed{
\ell_i
=
\max\{0,\mu_i-\varepsilon\}
\le
\lambda_i
\le
u_i
=
\max\{0,\mu_i+\varepsilon\}.
}
\]
\end{corollary}

This remains valid even if symmetric phase pruning makes \(H_0\) slightly
indefinite.

\subsection{Certified generalized-\texorpdfstring{\(\chi^2\)}{chi-square} tail bracket}

For independent standard normal variables \(Z_i\), define
\[
Q_-=\sum_i\ell_iZ_i^2,
\qquad
Q=\sum_i\lambda_iZ_i^2,
\qquad
Q_+=\sum_i u_iZ_i^2.
\]

\begin{theorem}[Stochastic tail enclosure]
For every threshold \(q\),
\[
\boxed{
\Pr(Q_-\ge q)
\le
\Pr(Q\ge q)
\le
\Pr(Q_+\ge q).
}
\]
\end{theorem}

\begin{proof}
Couple all three quadratic forms using the same normal variables \(Z_i\).
Since
\[
0\le\ell_i\le\lambda_i\le u_i
\]
for every \(i\),
\[
Q_-\le Q\le Q_+
\]
almost surely.  Taking upper-tail events yields the result.
\end{proof}

The two endpoint probabilities can be evaluated by the same K-COS machinery
used for the unpruned spectrum.  Its numerical error certificate is separate
from the spectral enclosure above.

\subsection{A dual-budget pruning rule}

The local Krylov certificate from the previous section and the global spectral
certificate answer different questions.  The former asks whether omitted
couplings matter to the current Krylov query; the latter asks whether they can
move the spectrum enough to matter to the final test.

A production schedule should therefore enforce
\[
\boxed{
\sum_{e\in\mathcal O}\beta_e^{(k)}
\le
\tau_K,
\qquad
\|\mathcal C_{\mathcal O}\|_2
\le
\tau_\lambda.
}
\]

The implementation uses a conservative greedy rule: phase pairs are considered
in increasing query-risk order, and a pair is omitted only when adding it
preserves both inequalities.  The rule is certified but is not claimed to be
globally optimal under the two simultaneous constraints.

\begin{figure}[H]
\centering
\includegraphics[width=.84\linewidth]{figure_14_dual_budget_tail_bridge.png}
\caption{Dual-budget phase pruning.  The query-local budget controls Krylov
propagation while the operator budget controls final spectral and tail
uncertainty.  A loose operator budget can admit a strong coupling and widen
the tail interval by several decades even when the Krylov-local budget still
passes.}
\label{fig:dual-budget-tail}
\end{figure}

In the deterministic ten-phase experiment, fixing the query budget at
\(0.05\) and the operator budget at \(0.005\) removed \(23\) of \(45\)
undirected phase pairs.  At a true probability of \(10^{-8}\), the certified
tail interval was
\[
10^{-8.11143}
\le
p
\le
10^{-7.89036},
\]
a width of about \(0.221\) decades.  With an operator budget of \(0.0075\),
\(34\) pairs were removed and the interval widened to about \(0.320\)
decades.  By contrast, allowing a strong coupling to be removed under a very
loose operator budget produced a multi-decade interval.

\subsection{Certify after prune}

There is a second route that does not use the pruned operator as the final
spectral approximation.  Let \(U\) be any orthonormal sample-space basis,
possibly generated cheaply from a pruned or approximate Krylov process, and
define
\[
B=U^\top A,
\qquad
F=(I-UU^\top)A,
\qquad
\delta=\|F\|_2.
\]

\begin{theorem}[Range-core singular-value certificate]
Let the singular values of \(A\) be \(\sigma_i(A)\), and let
\(\widetilde\sigma_i(B)\) denote the singular values of \(B\) padded by zeros
to the ambient spectral dimension.  Then
\[
\boxed{
\left|
\sigma_i(A)
-
\widetilde\sigma_i(B)
\right|
\le
\delta.
}
\]
Consequently,
\[
\max\{0,\widetilde\sigma_i(B)-\delta\}^2
\le
\lambda_i(AA^\top)
\le
\left(
\widetilde\sigma_i(B)+\delta
\right)^2.
\]
\end{theorem}

\begin{proof}
Because \(U\) has orthonormal columns,
\[
A=UB+F
\]
and
\[
\|A-UB\|_2=\|F\|_2=\delta.
\]
The nonzero singular values of \(UB\) are exactly those of \(B\).
The standard singular-value perturbation inequality gives the stated bound.
Squaring the nonnegative endpoints gives the eigenvalue interval.
\end{proof}

This provides a useful architecture:
\[
\boxed{
\text{pruned operator for basis discovery}
\to
\text{one exact streamed factor check}
\to
\text{certified spectral core}.
}
\]

\subsection{A retained negative result}

The preceding theorem is exact, but a correct interval need not be useful.

In a full-rank \(50\times50\) stress test, a pruned Krylov basis of dimensions
\(12\) through \(48\) left range residuals between roughly \(0.87\) and
\(0.99\) times the leading singular value.  The corresponding \(10^{-8}\)
tail intervals were therefore vacuous or thousands of decades wide.  At
dimension \(50\), the range residual fell to approximately
\(1.2\times10^{-15}\), and the tail interval collapsed to the exact value.

\begin{figure}[H]
\centering
\includegraphics[width=.84\linewidth]{figure_15_range_core_certification.png}
\caption{Certify-after-prune range residual.  The range-core theorem is valid
at every depth, but a deep-tail full-spectrum guarantee can remain useless
until the range residual is genuinely small.  This negative result is retained
as a stopping-rule lesson rather than hidden as an implementation failure.}
\label{fig:range-core-cert}
\end{figure}

The practical lesson is that range-core certification is a stopping test, not
a promise that a severely incomplete basis automatically yields a useful
full-spectrum tail bound.

